% =============================================================================
% Analyse-Pipeline: saubere vertikale Darstellung
% =============================================================================
\begin{figure}[p]
  \centering
  \begin{tikzpicture}[
    scale=0.82, transform shape,
    % Styles
    step/.style={
      draw, rounded corners=4pt, fill=#1,
      text width=12cm, align=left, minimum height=1cm,
      font=\small, inner sep=8pt
    },
    step/.default={blue!6},
    filebox/.style={
      draw, fill=yellow!12, rounded corners=2pt,
      font=\scriptsize\ttfamily, inner sep=4pt
    },
    heading/.style={
      font=\small\bfseries, anchor=west
    },
    arr/.style={-{Stealth[length=5pt]}, thick, gray!70},
    loopback/.style={-{Stealth[length=5pt]}, thick, red!50!black, rounded corners=4pt},
  ]

    % ── 1. Messung ──────────────────────────────────────────────────────────
    \node[step=blue!8] (S1) at (0, 0) {
      \textbf{1 \enspace Synchrotron-Messung}\\[3pt]
      Absorptionsspektrum $\mu(E)$ im Bereich
      \num{7500}--\SI{10000}{eV}.
      Ein einziger Energiescan deckt die K-Kanten
      von Co (\SI{7709}{eV}), Ni (\SI{8333}{eV}) und
      Cu (\SI{8979}{eV}) gleichzeitig ab.
    };

    % ── 2. Athena ───────────────────────────────────────────────────────────
    \node[step=orange!8, below=0.45cm of S1] (S2) {
      \textbf{2 \enspace Vorverarbeitung in Athena}\\[3pt]
      Import der Rohdaten, visuelle Kontrolle,
      erste $E_0$-Bestimmung.
      Export als Athena-Projekt
      \,\tikz[baseline=-0.3ex]{\node[filebox]{MEA\_athena.prj};}
    };
    \draw[arr] (S1) -- (S2);

    % ── 3. Extraktion ──────────────────────────────────────────────────────
    \node[step=orange!8, below=0.45cm of S2] (S3) {
      \textbf{3 \enspace $\chi(k)$-Extraktion mit larch}\\[3pt]
      Neu-Extraktion nötig, da Multi-Edge-Normierung in
      Athena fehlerhafte Kantensprünge lieferte.
      Pro Kante wird ein separates Energiefenster gesetzt
      (Co: ${\leq}\SI{8233}{eV}$,
       Ni: \SIrange{8100}{8879}{eV},
       Cu: ${\geq}\SI{8700}{eV}$).\\[3pt]
      Schritte: $E_0$ bestimmen $\to$
      Pre-Edge abziehen $\to$
      auf $\Delta\mu_0 = 1$ normieren $\to$
      Hintergrund entfernen (\texttt{autobk}, $R_\text{bkg}=\SI{1.0}{\AA}$)
      $\to$ Umrechnung $E \to k$.\\[4pt]
      Ausgabe:\enspace
      \tikz[baseline=-0.3ex]{\node[filebox]{MEA\_Co\_chi.dat};}\enspace
      \tikz[baseline=-0.3ex]{\node[filebox]{MEA\_Ni\_chi.dat};}\enspace
      \tikz[baseline=-0.3ex]{\node[filebox]{MEA\_Cu\_chi.dat};}\\
      {\scriptsize Jeweils zwei Spalten: $k$ (\AA$^{-1}$) und $\chi(k)$.}
    };
    \draw[arr] (S2) -- (S3);

    % ── 4. Startmodell ─────────────────────────────────────────────────────
    \node[step=green!8, below=0.45cm of S3] (S4) {
      \textbf{4 \enspace Startmodell (Superzelle)}\\[3pt]
      Kristallstruktur wählen (FCC, BCC oder HCP) und
      $4{\times}4{\times}4$-Superzelle mit 256 Atomen
      (85\,Co, 85\,Ni, 86\,Cu, zufällig verteilt) erzeugen.
      Format: POSCAR mit Gittervektoren und fraktionellen
      Koordinaten.\\[4pt]
      Ausgabe:\enspace
      \tikz[baseline=-0.3ex]{\node[filebox]{MEA\_CoNiCu.p1};}
    };
    \draw[arr] (S3) -- (S4);

    % ── 5. Konfiguration ──────────────────────────────────────────────────
    \node[step=green!8, below=0.45cm of S4] (S5) {
      \textbf{5 \enspace Konfiguration}
      \hfill\tikz[baseline=-0.3ex]{\node[filebox]{parameters.dat};}\\[3pt]
      Zentrale Steuerdatei mit allen Einstellungen:
      $k$-Bereich (\SIrange{3}{11.5}{\per\AA}),
      $k$-Gewichtung ($k^2$),
      Vergleichsraum (Wavelet),
      Streupfade ($N_\text{legs}=4$),
      FEFF-Clusterradius (\SI{6}{\AA}),
      Schrittweite, Konvergenzkriterien.\\[3pt]
      Die Option \texttt{-autoES} bestimmt $S_0^2$ und $\Delta E_0$
      automatisch vor dem Hauptlauf.
    };
    \draw[arr] (S4) -- (S5);

    % ── 6. EvAX-Kern ──────────────────────────────────────────────────────
    \node[step=red!6, below=0.45cm of S5,
          draw=red!40!black, thick] (S6) {
      \textbf{6 \enspace EvAX-Optimierung}
      \hfill{\small\texttt{./EvAX.exe parameters.dat -autoES}}\\[4pt]
      %
      \begin{minipage}[t]{5.5cm}
        \textbf{Iterationsschleife:}\\[2pt]
        {\scriptsize
        \textbf{a)} Zufälliges Atom verschieben\\
        \quad\enspace (max.\ \SI{0.005}{\AA}/Schritt)\\[1pt]
        \textbf{b)} FEFF berechnet Streuphasen\\
        \quad\enspace für die neue Umgebung\\[1pt]
        \textbf{c)} $\chi_\text{calc}(k)$ mitteln über alle\\
        \quad\enspace Absorber gleichen Typs\\[1pt]
        \textbf{d)} Vergleich mit $\chi_\text{exp}(k)$\\
        \quad\enspace im Wavelet-Raum\\[1pt]
        \textbf{e)} Metropolis-Akzeptanz:\\
        \quad\enspace besser $\to$ behalten,\\
        \quad\enspace schlechter $\to$ evtl.\ behalten\\[1pt]
        $\to$ zurück zu \textbf{a)}
        }
      \end{minipage}%
      \hfill
      \begin{minipage}[t]{5.5cm}
        \textbf{FEFF-Dateien im Arbeitsverz.:}\\[2pt]
        {\scriptsize
        \tikz[baseline=-0.3ex]{\node[filebox]{feff};} Binary (feff85L)\\[2pt]
        \tikz[baseline=-0.3ex]{\node[filebox]{run-FEFF.exe};} Shell-Wrapper\\[2pt]
        \tikz[baseline=-0.3ex]{\node[filebox]{feff\_0/ feff\_1/};} je Kante:\\
        \quad \texttt{feff.inp} -- Eingabe (Cluster)\\
        \quad \texttt{feff.bin} -- Streuphasen\\
        \quad \texttt{chi.dat} -- berechnetes $\chi(k)$\\[4pt]
        \textbf{Evolutionärer Algorithmus:}\\[2pt]
        16 parallele Strukturkopien (States).\\
        Alle $\sim$100 Iterationen: schlechte\\
        States durch gute ersetzen +\\
        leichte Mutationen.
        }
      \end{minipage}
    };
    \draw[arr] (S5) -- (S6);

    % ── 7. Ausgaben ───────────────────────────────────────────────────────
    \node[step=purple!6, below=0.45cm of S6] (S7) {
      \textbf{7 \enspace Ausgaben}
      \hfill{\scriptsize Verzeichnis \texttt{output\_files/}}\\[4pt]
      \tikz[baseline=-0.3ex]{\node[filebox]{output.dat};}
        {\scriptsize\enspace Residual vs.\ Iteration (Konvergenz)}\hfill
      \tikz[baseline=-0.3ex]{\node[filebox]{final.xyz};}
        {\scriptsize\enspace optimierte 256 Atompositionen}\\[3pt]
      \tikz[baseline=-0.3ex]{\node[filebox]{Co/EXAFS.dat};}
        {\scriptsize\enspace $\chi_\text{calc}(k)$ und $\chi_\text{exp}(k)$
         pro Kante}\hfill
      \tikz[baseline=-0.3ex]{\node[filebox]{Co/expFT.dat};}
        {\scriptsize\enspace Fouriertransformierte}\\[3pt]
      \tikz[baseline=-0.3ex]{\node[filebox]{Co/RDF\_Co\_Ni.dat};}
        {\scriptsize\enspace partielle Paarverteilungsfunktionen $g_{ij}(R)$}\hfill
      \tikz[baseline=-0.3ex]{\node[filebox]{restart*};}
        {\scriptsize\enspace Checkpoint zum Fortsetzen}
    };
    \draw[arr] (S6) -- (S7);

    % ── 8. Auswertung ────────────────────────────────────────────────────
    \node[step=green!10, below=0.45cm of S7,
          draw=green!40!black] (S8) {
      \textbf{8 \enspace Auswertung (Python)}\\[4pt]
      \begin{minipage}[t]{5.5cm}
        \tikz[baseline=-0.3ex]{\node[filebox]{plot\_results.py};}\\[2pt]
        {\scriptsize Parst \texttt{output\_files/},
         erzeugt Konvergenz-, EXAFS-,
         FT-, RDF- und SRO-Plots.}
      \end{minipage}%
      \hfill
      \begin{minipage}[t]{5.5cm}
        \tikz[baseline=-0.3ex]{\node[filebox]{parameter\_scan.py};}\\[2pt]
        {\scriptsize Automatisierte Variation aller
         Fit-Parameter (83 Kombinationen,
         6 parallele EvAX-Instanzen).}
      \end{minipage}\\[6pt]
      \textbf{Physikalische Ergebnisse:}\enspace
      {\scriptsize
      Koordinationszahlen $N_{ij}$,\enspace
      Bindungslängen $\langle R_{ij}\rangle$,\enspace
      Warren-Cowley-Parameter $\alpha_{ij}$,\enspace
      Kristallstruktur.}
    };
    \draw[arr] (S7) -- (S8);

  \end{tikzpicture}
  \caption{Vollständige Analyse-Pipeline von der Synchrotron-Messung
    bis zu den physikalischen Ergebnissen.
    Gelbe Kästen bezeichnen konkrete Dateien.
    Der rot umrandete Schritt~6 zeigt den iterativen Kern von \evax{},
    in dem FEFF-Rechnungen und RMC-Optimierung zusammenwirken.}
  \label{fig:full-pipeline}
\end{figure}
